% ── CHAPTER 8: CONVERGENCE 2: NONLOCAL INTERCONNECTION (~6,000 words) ──
\chapter[Nonlocal Interconnection]{Nonlocal Interconnection: The Bond That Ignores Distance}
\label{ch:nonlocal-interconnection}

% \epigraph{Spooky action at a distance.}{Albert Einstein, letter to Max Born, 1947}

\firstword{I}{n} 1935, Albert Einstein co-authored a paper designed to prove that quantum mechanics was incomplete \citep{Einstein1935EPR}. The argument, later called the EPR argument after its authors (Einstein, Podolsky, and Rosen), rested on a thought experiment. Imagine two particles that interact and then fly apart in opposite directions. Quantum mechanics predicts that measuring the spin of one particle will instantly determine what result you will get if you measure the spin of the other, no matter how far apart the particles are. (The original paper used position and momentum; the version with spin, which is easier to picture and closer to what was later tested, is due to David Bohm.) Einstein found this intolerable. If no signal can travel faster than light --- and Einstein had spent his career establishing this --- then the second particle's measurement outcome cannot be caused by what happens to the first. Therefore, Einstein argued, the outcomes must have been determined in advance. The particles must carry ``hidden variables'' --- pre-set instructions --- that tell them what outcome to produce when measured. Quantum mechanics, by failing to include these variables, was incomplete.

Einstein was right to be disturbed, and right that something is going on that has no classical precedent. Where he was wrong --- and where the situation becomes far more interesting than he perhaps wanted it to be --- was in the conclusion. The hidden variables do not save him. In 1964, John Bell proved mathematically that any theory with the kind of pre-existing local hidden variables Einstein was imagining would have to satisfy certain statistical constraints \citep{Bell1964}. Experiments by Alain Aspect and his colleagues in 1982 \citep{Aspect1982}, and the more stringent tests that followed, gave the answer. The constraints are violated. The correlations between entangled particles are stronger than any pre-set instructions could produce. Einstein's intuition that something is going on was correct. His proposed remedy was not.

What is going on is entanglement. And the thing to understand about entanglement is that it is not a weak or metaphorical kind of connection. It is a genuine feature of physical reality --- confirmed by half a century of experiments, honored with a Nobel Prize in 2022 \citep{Zeilinger2023Nobel}, and now put to practical use in quantum cryptography and quantum computing.

The second convergence is about this feature. Three traditions of inquiry --- quantum physics, process philosophy, and the \bahai writings --- each describe, in their own terms, a kind of bond that persists between things regardless of spatial separation. And each describes this bond in the same structural way: not as a signal sent from one place to another, nor simply as a causal influence transmitted through space, but as a fact about the constitution of reality itself --- a relational unity that is not destroyed by distance because it was never a matter of distance in the first place.

That is the convergence I want to demonstrate in this chapter. To do it properly, I need to begin with what entanglement actually is: not the popular version, which is often oversimplified or mystified, but the structure the formalism describes and the experiments confirm.

% ─────────────────────────────────────────────────────────────────────────────
\section{The Physics: Entanglement and the Nonlocal Bond}
\label{sec:ch8-physics}
% ─────────────────────────────────────────────────────────────────────────────

Start with what entanglement is, stated carefully.

Two particles are entangled when their joint quantum state cannot be written as the product of two separate individual states. This sounds like a technical point, but its implications are enormous. It means the two particles are not, ontologically speaking, two separate things that happen to be correlated. They are one composite system, and that composite system has a complete quantum description, what physicists call a pure state, even though neither particle, considered alone, has one. Chapter~\ref{ch:relational-ontology} worked through the simplest case, a pair of electrons whose total spin is zero while neither electron has a definite spin direction of its own. Here I want to ask what such a pair does when its two members are far apart.

An analogy that is sometimes used, though it is imperfect: imagine two gloves. If you send one glove to Paris and one to Tokyo, then the moment you open the Paris box and discover a left glove, you instantly know the Tokyo glove is right. But this is a classical correlation, not entanglement. The gloves had their handedness from the beginning; the measurement just revealed it. Einstein's hidden-variable intuition was, essentially, that entangled particles are like the gloves.

Bell's theorem proves they are not. The quantum correlations between entangled particles are not of the glove type. They cannot be explained by any pre-set property the particles carry with them, so long as what happens at one detector cannot depend on what is done at the other. The most vivid way to see this: if you measure both particles in the same direction, the quantum correlations are perfect. But if you rotate the measurement axes for the two particles by various angles, the patterns of correlation and anti-correlation follow a curve that no pre-existing set of instructions could reproduce. The pattern depends on the angle between the two settings, and no list of answers fixed at the source can match that dependence for every pair of settings the experimenters might choose. The worked example at the end of this section shows why, with nothing more than counting.

The key clarification --- and it matters enormously for the philosophical discussion --- is that this nonlocal connection does not allow faster-than-light communication. The point is misunderstood in both directions. Some conclude that a connection useless for communication cannot be real connection; others that quantum nonlocality lets us influence things instantly across any distance. Neither conclusion is right. When you measure an entangled particle, you get a random outcome that you do not choose. What you choose is which direction to measure, and that choice affects the correlations, but only in a way that becomes apparent when the two sets of results are compared afterward, which requires ordinary, slower-than-light communication. The nonlocal connection is real, but it is not a communication channel.

What, then, is it?

Quantum information theory treats entanglement as a resource: something the universe has that classical physics lacks, a nonlocal correlation that can be quantified, manipulated, and put to work. The physicist William Wootters, who has worked at the foundations of that field for decades, gave an exact formula for how much entanglement, on one standard measure, any state of two qubits contains \citep{Wootters1998}. This perspective is enormously useful for practical purposes. But for philosophical purposes I want to press a bit deeper.

What the formalism is telling us, in the most ontologically direct interpretation, is that entangled particles constitute a single composite system whose unity is not diminished by spatial separation. This is an ontological claim, not merely a claim about what information we can extract. The individual particles, considered alone, do not have complete states of their own. The composite system does. The unity is real; the individuality is what is incomplete or derivative.

Classical physics assumed that the world is made up of locally separate things that interact when they are close together and become independent when they are far apart. Distance severs connection. For most everyday purposes that picture is exactly right.

Entanglement tells us this picture is wrong at the fundamental level. Once two quantum systems have interacted and become entangled, they remain a single composite system however far apart they travel, so long as nothing else disturbs them. The bond is not transmitted through space; it is not a field that weakens with distance. It is a fact about the joint quantum state, and distance alone does nothing to that state. What does erode it is further interaction with other systems, which spreads the entanglement into the surroundings; that process, decoherence, is the subject of Chapter~\ref{ch:decoherence-decomposition}. Measuring one system does not send a signal to the other; it resolves the state of the composite system, and that resolution is reflected in what the other system will show if measured.

The popular analogy that I find most helpful for conveying the ontological point is this: entangled particles are not like twins who telephone each other when something important happens. They are like twins who share a single dream. In the first case, two separate people have their own experiences, and one sends news to the other across the distance between them. In the second, there is one dream, and both twins are in it. Separating the twins does not split the dream, because the dream is not located in either twin.

This analogy should not be taken too literally; dreams are not quantum states, and the physics is precise where analogies are not. But it captures the ontological point: the unity of entangled particles is the oneness of a composite system, and spatial separation is simply not a variable in the definition of that oneness.

\subsection*{A Worked Example: A Game That Instructions Cannot Win}

The heart of Bell's argument needs nothing more than counting. I will use the form that John Clauser, Michael Horne, Abner Shimony, and Richard Holt made testable in 1969 \citep{CHSH1969}, set out as a game.

Two players, Alice and Bob, are placed far apart and cannot communicate. In each round a referee tosses a coin for each player and asks that player question 0 or question 1 accordingly. Each answers $+1$ or $-1$. They win the round if their answers agree, except when both were asked question 1; then they win if their answers differ. Beforehand they may agree on any strategy they like.

Suppose they agree on instructions: for each question, the answer to give. To win every round in which at least one of them is asked question 0, Alice's answer to question 0 must match both of Bob's answers, and her answer to question 1 must match his answer to question 0. Chain these together and her answer to question 1 matches his answer to question 1, so they lose whenever both are asked question 1. Every set of instructions therefore fails on at least one of the four combinations of questions, each of which comes up a quarter of the time. Shuffling instruction sets at random, or making the answers probabilistic, only averages over strategies that each fail somewhere. The best any instructions can do is to win three rounds in four: 75 percent. In the standard notation, a certain sum of four correlations, written $S$, cannot exceed 2.

Now give the players entangled photon pairs, one photon of each pair to each, and let them answer by measuring polarization: each photon leaves the polarizer by one of two routes, recorded as $+1$ or $-1$. For a suitably prepared pair, the two answers agree with a probability equal to the square of the cosine of the angle between the two polarizers. Let Alice use 0 degrees for question 0 and 45 degrees for question 1, and Bob 22.5 degrees for question 0 and $-22.5$ degrees for question 1. In the three kinds of round where agreement wins, the polarizers are 22.5 degrees apart, and the answers agree about 85 percent of the time. In the fourth, they are 67.5 degrees apart, and the answers disagree about 85 percent of the time. In the ideal case the pairs win about 85 percent of all rounds; $S$ reaches $2\sqrt{2}$, about 2.83, and no quantum strategy can do better. Experiments built on this logic, from Aspect's in 1982 \citep{Aspect1982} to the tests of 2015 that closed the main loopholes together \citep{Zeilinger2023Nobel}, find violations of the classical ceiling, in agreement with quantum mechanics.

What the violation rules out is the combination of assumptions the instructions embody: that whatever settles each answer, as a definite value or only as a probability, was fixed at the source and travels with the photon; that each answer depends only on that and on its own polarizer; that the questions are chosen independently of what the photons carry; and, tacitly, that each round ends in one pair of results that is a fact for everyone. At least one must go, and the interpretations of quantum mechanics differ over which. Bohmian mechanics lets the result at one station depend on the setting at the other. The relational reading of Chapter~\ref{ch:relational-ontology} gives up the tacit assumption: each result is a fact relative to the station that recorded it, and their pairing is a fact only relative to something that has interacted with both. A few theorists deny that the settings are ever independent of the source. This is the sense in which physicists call the correlations of entangled pairs nonlocal: they cannot be built from contributions confined to each station plus whatever the pair carried from the source.

The same example shows why entanglement cannot carry a message. Whatever Bob does, whichever question he is asked or whether he measures at all, Alice's own list is a patternless run of $+1$s and $-1$s in equal numbers. Taken alone, each photon is in the kind of maximally mixed state described in Chapter~\ref{ch:relational-ontology}, and nothing done to its partner can change that \citep{NielsenChuang2000}. Bob's choice shows up only in the correlation between the two lists, which appears when the lists are brought together by some ordinary channel. When the two measurements are far enough apart that no light signal could pass between them, relativity adds that there is no fact about which came first: observers moving in different ways disagree about the order. A story in which Alice's measurement reaches across and fixes Bob's photon can be told just as well in reverse, with identical predictions. The correlation belongs to the pair. Distance does not appear in the calculation at all; only the angle between the polarizers does.

% ─────────────────────────────────────────────────────────────────────────────
\section{The Process-Philosophical View: Universal Prehension}
\label{sec:ch8-process}
% ─────────────────────────────────────────────────────────────────────────────

Whitehead's philosophy of nature was not constructed in response to quantum mechanics, though Whitehead was aware of the early quantum theory. It was constructed in response to a more general problem: how to understand a universe that, on every scientific account available by the early twentieth century, turned out to be not a collection of locally independent substances but a web of interconnected processes.

The aspect of Whitehead's framework that is directly relevant here is his account of prehension as universal. In the previous chapter, I described prehension as the relational act through which each actual occasion grasps prior occasions and incorporates them into its own becoming. The point I want to develop now is the scope of this prehension: it is not limited by spatial proximity.

In Whitehead's formulation, every actual occasion prehends all prior occasions \citep{Whitehead1929}. Not just the spatially adjacent occasions. Not just the most recent occasions. The entire relevant causal past is, in principle, available for prehension by any new occasion as it comes into being. Whitehead states this with characteristic directness: every actual entity is present in every other actual entity, in the sense that the objective character of each past occasion is available to every subsequent occasion's becoming.

For a past occasion to be ``present in'' a later one means that the definite character it achieved, what it was and what relational facts it established, becomes part of the data out of which the new occasion constitutes itself. A past event on the other side of the universe is, in principle, part of the relational field from which any new occasion emerges.

Of course, in practice, the prehension of very distant past occasions will be extremely weak --- attenuated by the structure of the causal past, which is organized by spacetime geometry. What Whitehead means is not that every past occasion has equal weight in every subsequent occasion's constitution, but that no past occasion is in principle sealed off from prehension by spatial distance alone. Distance affects the intensity of prehension; it does not determine whether prehension is possible at all.

This is a significant structural claim, and I want to be exact about its scope. What Whitehead makes universal is inheritance from the past: an occasion may draw on any event in its causal past, however remote. That is not yet nonlocality in the physicist's sense, since everything in an event's causal past could in principle have sent it a signal. What it does mean is that, for Whitehead, distance is not what separates things.

Think of a river carrying sediment. A drop of water in the lower Mississippi is not constituted only by the geography of the last few miles; it carries the chemical signature of the Rockies, the Missouri River basin, the Tennessee watershed, the whole upstream history attenuated and mixed but genuinely there. Whitehead's prehension works similarly, except that the inheritance is of relational character rather than chemical traces. The events that came before, all of them in principle though at widely varying intensities, are prehended by the new event and incorporated into its becoming. Its ``location'' is a feature of its most intense prehensions, not a boundary of the prehension relation itself.

This framework was developed a generation before Bell's theorem and the experimental confirmation of quantum nonlocality. What it offers the physics is a picture of reality in which relational bonds persist independently of spatial separation, because the bonds are constitutive facts about the becoming of occasions rather than transmitted effects propagating through space.

In the Whiteheadian picture, when two particles interact --- when they become entangled, in the quantum mechanical language --- each prehends the other, and the prehensive relation is established as an objective fact. This fact does not subsequently travel anywhere, and it is not transmitted between the particles when one is measured. It is a standing feature of the relational constitution of both, present in the objective past from which any subsequent occasion draws its inheritance. When the entanglement is tested, the result is not caused by the measurement event traveling across space to influence the other particle. Rather, both measurement events draw on the same relational heritage --- the objective fact of the past interaction that constituted both particles as aspects of a single relational whole.

That heritage has to be understood in a particular way, or the account slides back into the gloves. If what each measurement inherited were a settled answer, the shared heritage would be a set of instructions, and Bell's theorem rules those out. The heritage has to function as what Whitehead calls real potentiality, and, I would add, as a potentiality that belongs to the pair as a whole and cannot be divided into one inheritance per particle. What that requires of Whitehead's scheme is a question I return to at the end of the chapter.

None of this means that Whitehead's philosophy predicts quantum entanglement, or that the quantum formalism can be derived from process principles; the empirical content of quantum mechanics requires the formalism. On this reading, quantum nonlocality is one place where the constitutive character of relations, which Whitehead put at the center of his scheme, shows itself in the formalism.

Michael Epperson has argued at length that Whitehead's categories can serve as an ontological interpretation of quantum mechanics \citep{Epperson2004}. My claim here is narrower: that process ontology's account of relational bonds and quantum mechanics' account of entanglement share a logical form. Both deny that the bond between two related things is something transmitted from one to the other through space, and both affirm that it is a constitutive fact that persists independently of spatial separation.

There is a passage in \textit{Process and Reality} that is worth quoting at length here because it expresses, in Whitehead's own terms, the ontological picture I am trying to convey:

\begin{quote}
``The extensive continuum is one relational complex in which all potential objectifications find their niche. It underlies the whole world, past, present, and future.''
\citep[pp.~79--80]{Whitehead1929} % TODO verify wording and page against the edition cited (Process and Reality is not in sources/)
\end{quote}

The extensive continuum is not space as a container; it is the relational field of all potential connections. Every actual occasion exists within this universal relational field, which is why no occasion is ever in principle disconnected from any other. The field is one; its nodes are many; and the many are constituted by their patterned positions within the one field. Spatial separation is a feature of the extensive continuum, not a barrier to the relational connections that the continuum makes possible.

% ─────────────────────────────────────────────────────────────────────────────
\section{The \bahai View: The Complete and Perfect Linkage}
\label{sec:ch8-bahai}
% ─────────────────────────────────────────────────────────────────────────────

Among all the passages in the \bahai writings that bear on the structural questions of this book, there is one that seems to me the most philosophically precise and the most directly relevant to the phenomenon of quantum nonlocality. It comes from a letter \AbdulBaha wrote in 1921, in the course of setting out proofs of the existence of God, and it is included in \emph{Selections from the Writings of \AbdulBaha}, a compilation of his letters:

\begin{quote}
``All created things are connected one to another by a linkage complete and perfect, even, for example, as are the members of the human body.''
\citep[\S 21]{SelectionsAbdul}
\end{quote}

I want to read this passage slowly and with philosophical care, because there are several features of it that matter for the convergence argument.

First, the scope: \emph{all created things}. Not some created things, not things in a particular spiritual domain, not human beings in their spiritual life. All created things. The claim is universal.

Second, the character of the linkage: \emph{complete and perfect}. Not approximate, not partial. This is a strong claim, and I do not think it should be softened in the interest of making it seem more modest. The passage says nothing explicit about distance. What it says is that the members of an endless universe are linked one to another, and that the linkage leaves nothing out and has no defect. I take that to exclude a linkage that binds near neighbors and fades beyond them, but the reading is mine, not the text's wording, and I will ask later in the chapter how much weight it can bear.

Third, the analogy used to clarify the nature of this linkage: the members of the human body. This analogy is philosophically precise, and I want to explain why.

The members of a body are causally connected, but they are also \emph{constitutively unified}. The heart-lung system is a single physiological unit in which the health of each part is intrinsic to the health of the whole, and the health of the whole to the health of each part. What happens to the heart later causes effects in the lungs, but it is first a change in the state of the whole organism, of which the lungs are an intrinsic part.

This gives two ways of reading the linkage. The weaker reading says that all created things influence one another at a distance, the way the sun's gravity influences the earth from ninety-three million miles away: causal connection, real but transmitted through the structure of a field. The stronger reading, which I think the body analogy points to, says that the linkage is constitutive as well as causal. The things connected are not prior independent substances that later affect one another; they are members of a single reality whose unity is inherent rather than transmitted, of the same logical type as the unity of a body and not merely that of planets drawn together by gravity. That is a claim about ontological structure, not about the mechanism of influence.

\AbdulBaha returns to this theme in a number of different contexts in his talks and writings. In a talk at All-Souls Church in Chicago on 5 May 1912, collected in \emph{The Promulgation of Universal Peace}, he states the principle that underwrites the universality of this linkage:

\begin{quote}
``Every contingent and phenomenal being is a composition of distinct elements. As long as there is affinity and cohesion among these constituent elements, strength and life are manifest; but when dissension and repulsion arise among them, disintegration follows.''
\citep[Talk 41]{PUP}
\end{quote}

Here the ontological dimension is explicit. The bond among the elements is the condition of a being's existing as an integrated whole, not an optional addition to a being that would exist anyway. \AbdulBaha draws a social conclusion from it in the next breath, that peace and amity are the saving factors of human society, but the principle he appeals to is stated about every composite thing: dissension does not merely weaken the whole; it dissolves it. This parallels the claim from the previous chapter that without affinity there is no integrated being.

The universal scope of the \bahai claim also resonates with a passage from \Bahaullah himself, about the nature of the relationship between all things and the divine ground of existence:

\begin{quote}
``Every created thing in the whole universe is but a door leading into His knowledge, a sign of His sovereignty, a revelation of His names.''
\citep[Gleanings LXXXII]{GleaningsBahaullah} % TODO verify Gleanings section number
\end{quote}

This is not directly a statement about nonlocal connection between created things. It is a statement about the relational structure of created existence: every created thing is constituted by its relationship to the creative ground, and through that shared relation all things are connected. On this reading the linkage is complete and perfect because the ground from which all things emerge is one.

I want to be careful here not to reduce the theological content of these passages to a physics-flavored paraphrase. The \bahai texts make claims about divine reality, spiritual development, and the nature of creation that go far beyond what I am using them to establish. My claim is only that, within that larger theological and spiritual context, there is a structural claim about the nature of existence --- about the constitutive character of relational bonds, and about the universality of those bonds --- that is the same, at the level of logical form, as the claim that quantum mechanics and process philosophy are making.

% ─────────────────────────────────────────────────────────────────────────────
\section{One Phenomenon, Three Descriptions}
\label{sec:ch8-convergence}
% ─────────────────────────────────────────────────────────────────────────────

Let me now make the structural isomorphism precise.

The popular version of this convergence would say something like: ``physics shows that things are connected, and religions have always said the same thing.'' That is an analogy. It works by pointing to a shared theme --- connectedness --- without specifying the logical structure of the claim being made. An analogy at this level of generality can be found between almost any physical phenomenon and almost any spiritual tradition. It proves nothing and illuminates little.

The structural isomorphism I am arguing for is more specific and more demanding.

All three traditions make a claim about nonlocal connection that has the following structure:

\begin{enumerate}
\item Two or more entities that have been in relational contact share a bond.
\item This bond persists regardless of spatial separation.
\item The bond is ontological, not merely epistemic: it is a fact about what the entities are, not merely about what we know about them.
\item The bond is constitutive: it is not reducible to a signal sent between the entities or an influence transmitted through space, but is a fact about the relational constitution of the entities themselves.
\item The individual states of the separated entities are not complete: each entity's full description requires reference to the bond.
\end{enumerate}

Check each tradition against this structure.

\emph{Quantum mechanics}: Entangled particles share a joint quantum state that persists regardless of spatial separation \citep{Bell1964,Aspect1982,Zeilinger2023Nobel}. The bond is ontological: the individual particles do not have complete states of their own; the composite system does. The bond is not transmitted through space; no signal passes between the particles. The individual descriptions are incomplete: a full description of either particle requires reference to the entangled pair state \citep{NielsenChuang2000}. Points 1 through 5: check.

\emph{Process philosophy}: Every actual occasion prehends all prior occasions, establishing relational bonds that persist in the objective past regardless of subsequent spatial separation \citep{Whitehead1929}. The bonds are ontological: they are constitutive facts about the occasions' character, not merely facts about what observers know. The bond between two separated occasions is a standing feature of their relational constitution rather than a transmission from one to the other. The individual occasions' descriptions are incomplete without reference to their prehensive inheritance. Points 1 through 5: check, with a qualification about how correlated outcomes are settled that I take up below.

\emph{\bahai writings}: All created things are connected by a linkage ``complete and perfect'' \citep[\S 21]{SelectionsAbdul}, of the same type as the constitutive unity of the members of a body. The linkage is ontological: it is a claim about the structure of created existence, not merely about how things seem. It includes influence, but what makes it this kind of bond is the constitutive unity of a whole. The individual things' descriptions are incomplete without reference to the whole of which they are aspects. Points 1 through 5: check.

What makes this a structural isomorphism rather than an analogy is that the three traditions share this five-point structure, not merely a theme. A conversation and a chemical bond both involve connection without being structurally isomorphic. Different vocabularies; identical logical form.

This convergence is worth examining from the perspective of how each tradition arrived at it.

Quantum mechanics arrived here through experiment and mathematics. Bell's theorem is a mathematical proof. Aspect's experiments are laboratory results. No prior philosophical commitment to nonlocality drove the physics to this conclusion; the physics was forced there by the evidence. The nonlocal bond was discovered, not chosen.

Process philosophy arrived here through conceptual analysis. Whitehead was analyzing what it would mean for an event of becoming to be genuinely constituted by its relational inheritance from the past. The universality of prehension follows from the nature of the analysis: if becoming is constituted by relational inheritance, and if the causal past is structured by the entire history of the universe and not merely by spatial adjacency, then prehension is universal. The bond was derived, not discovered empirically.

The \bahai writings arrived here through a different path again: a theology of divine unity working out its implications for created existence. If all things are created by and through one divine reality, and their diversity is a diversity of relational expressions of that reality rather than a plurality of independent substances, then all things share a constitutive bond that is not a matter of spatial proximity. (The 1921 letter runs the argument in the other direction, from the linkage of all things to a single Power that directs them.) The bond was articulated from within a tradition of revealed and interpreted spiritual understanding.

Three routes. Three vocabularies. One structure. This is the second convergence.

\subsection*{The Strongest Objection, and What the Convergence Does Not Claim}

The strongest objection turns this chapter's physics against the rest of it. I stressed that entanglement is not a channel of influence: nothing Bob does makes any difference Alice can detect. The 1921 letter describes its linkage in quite different terms. The members of the body act on one another, so that ``the eye must look ahead before the step is taken,'' and when one member fails the others suffer for it \citep[\S 21]{SelectionsAbdul}. A little further on, the letter describes how created things are subject to influences from outside themselves, the sun warming, the rain nourishing, rain depending on clouds and clouds on the sun, and argues from that chain to ``One Who influenceth all, and yet is influenced by none'' \citep[\S 21]{SelectionsAbdul}. Whitehead's prehension, likewise, is causal inheritance from the past. On this objection, physics describes a correlation that carries no influence, the other two traditions describe webs of influence, and my five-point structure hides the difference under the word ``bond.''

I think the objection is right about the letter and right about Whitehead, and the convergence has to be narrowed to meet it.

The letter's linkage plainly includes ordinary causation running through intervening links, and nothing in it suggests an influence that skips them. What the body analogy adds to that web is a claim about the whole: each member has its own special function, which it has only as a member, and a single directing power coordinates them all \citep[\S 21]{SelectionsAbdul}. That holistic element, and no particular kind of influence, is what I claim for the convergence.

The objection bites harder on Whitehead. In his scheme, contemporaries, occasions neither of which lies in the other's past, happen in causal independence of one another \citep{Whitehead1929,Whitehead1933}. The two measurements in a Bell experiment are contemporaries in exactly this sense, so neither prehends the other. If each simply settled its outcome from its own inheritance, the result would be a local picture of the kind the worked example rules out. The process reading survives only if the shared inheritance is a joint potentiality that the two contemporaries settle together. I think that is in the spirit of what Whitehead says about the extensive continuum, but he did not work it out, and I offer it as an extension of his scheme, not as his doctrine.

What survives is a common structure at the level of wholes and parts. In all three traditions the whole has a unity that is not assembled from separately complete parts; a part's full description refers to the whole; and that unity is not something sent across the distance between the parts, even where, as in the body and in the causal past, the parts also act on one another. The shared structure is holism. There is no shared mechanism, and I do not claim one.

With that narrowing in place, here is what the convergence does not claim.

\begin{itemize}
\item It does not claim that anything travels faster than light, or that entanglement could carry a message, an intention, or a healing influence from one person to another.
\item It does not claim that everything affects everything else in a causal sense. The causal web of the 1921 letter runs through intervening links, and on no interpretation can entanglement be used as a causal link between two places.
\item It does not claim that every created thing is entangled with every other. Interaction generically produces entanglement, but further interaction spreads it into the surroundings rather than leaving it between two particular things, and between warm, constantly interacting bodies it is dispersed almost at once (Chapter~\ref{ch:decoherence-decomposition}). People who love one another are not entangled in the physicist's sense.
\item It does not claim to settle which of Bell's assumptions fails.
\item It does not claim that the \bahai writings describe entanglement, or that Whitehead foresaw Bell's theorem.
\end{itemize}

\begin{convergencemoment}{Nonlocal Interconnection: The Bond That Ignores Distance}
\textbf{From the quantum side.} Two particles that have interacted can share a single joint quantum state that persists however far apart they travel \citep{Bell1964,Aspect1982,Zeilinger2023Nobel}. The pair has a complete description; neither particle has one of its own. Their correlations cannot be built from anything each particle carries separately, yet they cannot be used to send a signal, and distance plays no part in them.

\textbf{From the process-philosophical side.} Every actual occasion inherits from its entire causal past, not merely from its spatial neighborhood \citep{Whitehead1929}. Occasions that share a relational event remain connected through the objective fact of that event, a connection that distance may attenuate but does not sever.

\textbf{From the \bahai side.} ``All created things are connected one to another by a linkage complete and perfect, even, for example, as are the members of the human body'' \citep[\S 21]{SelectionsAbdul}. The body analogy specifies the \emph{type} of connection: the unity of members within one whole, which is more than causal influence transmitted at a distance.

\textbf{The structural isomorphism.} All three describe a constitutive relational bond that is not severed by spatial separation, that is ontologically real rather than merely epistemic, and that is not reducible to anything transmitted through space. The shared structure is holism, and no shared mechanism is implied: nothing here travels faster than light, and nothing makes everything a cause of everything else. \emph{The bond that connects related things is not a thread that stretches between them across space. It is a fact about what they jointly are.}
\end{convergencemoment}
